\Hefteintrag{2}{Strukturierung von Programmcode - Semantik: Methoden}{%

Beim Programmieren können eigene Anweisungen/ Befehle/ Methoden/ Funktionen \textbf{deklariert} werden. Bei der \textbf{Deklaration beschreibt} zunächst mit dem \LoesungLuecke{Methodenkopf}{6cm}, welche Eigenschaften die Methode hat und anschließend im \LoesungLuecke{Methodenrumpf}{6cm}, wie genau sie sich verhalten soll. Eine deklarierte Methode kann anschließend beliebig oft verwendet (=\emphColA{aufgerufen}) werden.

Hierdurch kann man sehr lange Programme in übersichtlichere Einzelteile zerlegen, sodass \LoesungLuecke{doppelter Code}{8cm}~reduziert wird, der Code durch aussagekräftige \LoesungLuecke{Methodennamen}{6cm}~besser verstanden werden kann und Fehler schneller gefunden werden.

In Java darf außerdem kein auszuführender Programmcode außerhalb von Methoden sein.

Die Deklaration besteht aus \emphColB{Methodenkopf} und \emphColC{Methodenrumpf} (von geschweiften Klammern umschlossen):

{\ttfamily
\emphColB{void tueEtwas()} \emphColC{\{}

~~~~\emphColC{karol.schritt(2);}
    
\emphColC{\}}
}

Der \emphColA{Aufruf} funktioniert wie ein einfacher Befehl: 

\emphColA{\texttt{tueEtwas();}}

}